\begin{abstract}
Trajectory data for traffic analytics, simulation calibration and infrastructure studies rarely record the vehicle type. We ask whether it can be read from motion alone: can 60-s windows of trajectory and CAN signals reveal the size class (three footprint classes) and the powertrain (battery-electric vs combustion, BEV vs ICE) of car models never seen in training? We release a pseudonymous dataset of 9,062 drives (about 3,950 h) from 1,431 vehicles and 174 car models, with held-out-car-model and held-out-device splits, window-, drive- and vehicle-level protocols and cluster-bootstrap intervals over car models. Four sequence architectures (ResNet, bi-GRU, Transformer, Mamba-2), hand-crafted statistics and physical features proposed by large language models (LLMs) serve as reference points; a pre-registered method study on model-disjoint internal folds adds car-model-balanced sampling, auxiliary physical targets, physics-safe augmentation, per-car-model loss re-weighting and, for BEV vs ICE, down-weighting of driving automation system (DAS) periods. On held-out car models, BEV vs ICE reaches 92.6--93.0 drive-level macro-F1 (95\% interval [87.1, 96.2] for the final bi-GRU); size improves from 62.9 to 67.1 at drive level but stays far from 90. Run unchanged on a public fleet of 4,061 devices absent from our library, the final models call the powertrain right for 86 of 88 car models and match the internal size result (63.2\% of devices against 63.5\% at vehicle level). Zero-training oracle analyses attribute the gap to label resolution and training-fleet composition rather than to evidence per drive; the architecture benchmark shows model capacity is not the bottleneck.
\end{abstract}
% W1: shorter title alternative, if wanted:
% "Reading Vehicle Size and Powertrain from Motion: A Trajectory--CAN Benchmark with Held-Out Car Models"
